\chapter{Giới thiệu đề tài}
\label{ch:gioithieu}

Các dịch vụ dựa trên vị trí (LBS) đã trở thành một phần không thể thiếu trong đời
sống số: điều hướng, tìm kiếm địa điểm, gọi xe, giao đồ ăn, quảng cáo cá nhân hóa
và các ứng dụng y tế. Sự tiện lợi này đến từ việc người dùng liên tục báo cáo vị
trí của mình cho nhà cung cấp dịch vụ. Tuy nhiên, dữ liệu quỹ đạo --- chuỗi vị trí
theo thời gian --- có thể tiết lộ những thông tin nhạy cảm bậc nhất về một cá
nhân: nơi ở, nơi làm việc, thói quen sinh hoạt, tình trạng sức khỏe, quan hệ xã
hội và niềm tin tôn giáo. Khi dữ liệu này bị lạm dụng, hậu quả trải từ theo dõi,
phân biệt đối xử cho tới đe dọa an toàn cá nhân.

Trong giai đoạn Thực tập 2, chúng tôi đã xây dựng một pipeline bảo vệ quỹ đạo
dựa trên Geo-Indistinguishability (Geo-I)~\cite{andres2013geo}: thêm nhiễu Laplace
phẳng vào từng điểm, giới hạn bán kính nhiễu trong ngưỡng QoS, loại bỏ các điểm
rơi vào tòa nhà/mặt nước, và làm khớp điểm giả vào lưới đường. Khảo sát văn liệu
sâu hơn (Chương~\ref{ch:lienquan}) cho thấy pipeline này tồn tại hai lỗi hình thức
nghiêm trọng: \emph{(i)} việc chặn bán kính nhiễu quanh vị trí thật làm hai vị trí
đủ xa nhau có support đầu ra rời nhau, khiến đảm bảo Geo-I thuần không còn đúng;
\emph{(ii)} việc lấy mẫu lại khi điểm giả rơi vào vùng cấm điều kiện hóa phân phối
đầu ra trên vị trí thật, đồng thời chính loại ``lọc theo tính hợp lý'' này là bề
mặt mà các tấn công tái tạo hiện đại như RAoPT~\cite{buchholz2022raopt} khai thác.

Luận văn này kế thừa mục tiêu của Thực tập 2 nhưng thay thế cơ chế lõi bằng bốn
cơ chế mới hoạt động \emph{ngay trên lưới đường}:

\begin{itemize}
  \item \textbf{REM} (Road-network Exponential Mechanism): cơ chế lũy thừa trên
  tập đỉnh của đồ thị đường OSM, thỏa $\eps$-Geo-I Euclid hình thức, và mọi vị trí
  công bố đều nằm trên đường \emph{theo cấu trúc} --- loại bỏ \emph{tín hiệu}
  off-road mà tấn công map-matching/RAoPT khai thác (mức kháng tấn công cần kiểm
  chứng trực tiếp).
  \item \textbf{T-REM}: REM bổ sung trọng số khả đạt (reachability) chỉ phụ thuộc
  vào \emph{điểm đã công bố trước đó} --- thông tin công khai --- nên đảm bảo
  riêng tư per-release không đổi, trong khi \emph{giảm} tín hiệu speed-implausibility
  mà tấn công liên kết vận tốc dựa vào (một soft regularizer, không phải hard
  reachable-set).
  \item \textbf{SM-REM}: T-REM bổ sung memoization theo lưới công khai --- cùng
  một chỗ luôn trả về cùng một điểm công bố --- chống tấn công averaging tại
  stay-point (khôi phục địa chỉ nhà) trong phạm vi một cache-lifetime, rủi ro
  thực tế nghiêm trọng mà các cơ chế nhiễu độc lập bỏ ngỏ.
  \item \textbf{PR-SM-REM}: thay quyết định memoize chính xác (vốn rò pattern
  revisit) bằng \emph{phép thử reuse có nhiễu} (kiểu predictive mechanism), cho
  chặn per-release $(\eps_{\text{test}}+\eps_{\text{release}})$-Geo-I hữu hạn ---
  đóng lỗ hổng tỷ-số-vô-hạn của memo chính xác, đổi một phần sức chống-averaging
  lấy một chặn hợp lệ. (Chặn w-event scalar đầy đủ cần một budget manager theo cửa
  sổ --- chưa cài, là hướng phát triển.)
\end{itemize}

Lưu ý xuyên suốt: các đảm bảo $\eps$-Geo-I được chứng minh cho \emph{ideal
real-arithmetic kernel}; hiện thực dấu-phẩy-động là một xấp xỉ số học (có
artefact hữu-hạn-độ-chính-xác, xem Mục~\ref{sec:prsmrem} và
Chương~\ref{ch:tongket}).

Đồng thời, đối chiếu với hệ thống thực tế dẫn tới việc \emph{phát biểu lại bài
toán}: LBS không cần vị trí chính xác mà cần mức granularity tối thiểu theo mục
đích; ràng buộc chất lượng vì thế được phát biểu dưới dạng $(\alpha,\delta)$-hữu
dụng ($\Pr[d(x,z)\le\alpha]\ge\delta$, khớp semantics trường \texttt{accuracy}
mà mọi ứng dụng đang tiêu thụ) thay vì một ngưỡng cứng $\delta$ quanh điểm thật
(vốn phá vỡ đảm bảo, Mệnh đề~\ref{prop:cap}); và giải thuật được đánh giá theo
threat model rút từ các tấn công đã xảy ra thật.

Đóng góp của luận văn tính đến bản cập nhật này:
\begin{enumerate}
  \item Phát biểu lại bài toán theo mô hình triển khai thực tế (local, on-device,
  granularity theo mục đích, bảo vệ endpoint) với bảng yêu cầu thiết kế ánh xạ
  vào các harm thực tế (Mục~\ref{sec:systemmodel});
  \item Phân tích hình thức các lỗi của pipeline Thực tập 2
  (Mục~\ref{sec:phantichbaseline});
  \item Bốn cơ chế REM, T-REM, SM-REM và PR-SM-REM kèm chứng minh đảm bảo riêng
  tư \emph{ở mức ideal real-arithmetic kernel} (Mục~\ref{sec:rem}, \ref{sec:trem},
  \ref{sec:smrem}, \ref{sec:prsmrem}; hiện thực dấu-phẩy-động là xấp xỉ hữu hạn
  độ chính xác, Nhận xét~\ref{rmk:finiteprec});
  \item Bộ khung đánh giá theo chuẩn văn liệu: một adversary suy luận
  \emph{dạng-REM emission} (chính xác cho REM, là \emph{proxy} cho các cơ chế
  khác --- dùng làm chẩn đoán, không xếp hạng riêng tư), tấn công theo dõi HMM,
  và tấn công averaging tại stay-point, cùng các độ đo tiện ích và tính thực tế
  (Chương~\ref{ch:thucnghiem});
  \item Thực nghiệm trên dữ liệu GPS thật (GeoLife, Bắc Kinh) trên đồ thị đường
  OSM thật, so sánh với Laplace phẳng và baseline Thực tập 2;
  \item Hệ thống mô phỏng trực quan ba góc nhìn (người dùng / máy chủ LBS / kẻ
  tấn công) minh họa các kịch bản sử dụng đời thực.
\end{enumerate}

Phần còn lại của luận văn được tổ chức như sau. Chương~\ref{ch:cosolythuyet}
trình bày cơ sở lý thuyết. Chương~\ref{ch:lienquan} khảo sát các công trình liên
quan và xác định khoảng trống nghiên cứu. Chương~\ref{ch:phuongphap} mô tả và
chứng minh các cơ chế đề xuất. Chương~\ref{ch:thucnghiem} trình bày thực nghiệm.
Chương~\ref{ch:tongket} tổng kết và nêu hướng phát triển.
